\subsection{Non-perturbative DBI Envelope}
  \label{subsec:nonperturbative-dbi}

  The weakly non-linear reduction derived in
  Section~\ref{subsec:weakly-nonlinear-dbi-reduction} provides the leading
  amplitude correction to the dispersion relation.
  However, the discrete Dirac--Born--Infeld dynamics allows the
  single-mode problem to be solved exactly on the $Q_8$ Cayley graph.

  Consider again a single eigenmode
  \[
    \chi_v(t) = q(t)\psi_n(v),
  \]
  with $\mathcal{L}\psi_n = \lambda_n \psi_n$ and $\|\psi_n\|_2=1$.
  Substituting this ansatz into the DBI Lagrangian
  Eq.~\eqref{eq:dbi-lagrangian} gives the exact one-mode Lagrangian
  \begin{equation}
    L_n[q]
    =
    -c_\chi^2
    \sum_{v\in Q_8}
    \sqrt{1-\frac{\psi_n(v)^2\dot q^{\,2}}{c_\chi^2}}
    -
    \frac12\lambda_n q^2 .
  \end{equation}

  \subsubsection*{Abelian sectors}

    For the nontrivial abelian representations $\rho_1,\rho_2,\rho_3$,
    the normalized eigenmodes satisfy
    \[
      |\psi_n(v)| = \frac{1}{\sqrt{8}}
      \qquad
      \forall v\in Q_8 .
    \]
    The Lagrangian therefore becomes
    \begin{equation}
      L
      =
      -8c_\chi^2
      \sqrt{1-\frac{\dot q^{\,2}}{8c_\chi^2}}
      -
      \frac12\lambda_n q^2 .
    \end{equation}

  \subsubsection*{Non-abelian sector}

    For the quaternionic sector $\rho_4$ in the canonical basis,
    the eigenmodes satisfy
    \[
      |\psi_n(v)| =
      \begin{cases}
        \frac12 & \text{on four vertices}, \\
        0 & \text{on four vertices}.
      \end{cases}
    \]
    The DBI sum therefore reduces to
    \begin{equation}
      L
      =
      -4c_\chi^2
      \sqrt{1-\frac{\dot q^{\,2}}{4c_\chi^2}}
      -4c_\chi^2
      -
      \frac12\lambda_n q^2 .
    \end{equation}
    The constant term does not affect the dynamics, and the effective
    Lagrangian becomes
    \begin{equation}
      L
      =
      -4c_\chi^2
      \sqrt{1-\frac{\dot q^{\,2}}{4c_\chi^2}}
      -
      \frac12\lambda_n q^2 .
    \end{equation}

  \subsubsection*{Energy conservation}

    Since the Lagrangian is time independent, the corresponding Hamiltonian
    is conserved.
    One obtains
    \begin{equation}
      E
      =
      4c_\chi^2
      \frac{1}{\sqrt{1-\frac{\dot q^{\,2}}{4c_\chi^2}}}
      +
      \frac12\lambda_n q^2 .
    \end{equation}
    At the turning points $\dot q=0$, where $q=\pm a_n$, the energy becomes
    \begin{equation}
      E
      =
      4c_\chi^2
      +
      \frac12\lambda_n a_n^2 ,
    \end{equation}
    where $a_n$ denotes the raw oscillation amplitude of the modal coordinate.

  \subsubsection*{Velocity bound}

    The DBI square root imposes the kinematic bound
    \begin{equation}
      |\dot q| < 2c_\chi .
    \end{equation}
    Expressed in terms of the field variable,
    \[
      \chi_v = q\psi_n(v),
    \]
    this corresponds exactly to the local flux constraint
    \[
      |\partial_t \chi_v| \le c_\chi .
    \]

  \subsubsection*{Amplitude bound}

    The turning point occurs when the maximal vertex saturates the flux bound,
    \[
      |\partial_t\chi_v|
      =
      |\dot q|\,|\psi_n(v)|
      =
      c_\chi .
    \]
    Since $\|\psi_n\|_\infty$ is the maximal value of $|\psi_n(v)|$, this gives
    \begin{equation}
      a_n^{\max}
      =
      \frac{c_\chi}{\|\psi_n\|_\infty\sqrt{\lambda_n}} .
    \end{equation}
    The corresponding effective amplitude
    \[
      A_n = a_n\|\psi_n\|_\infty
    \]
    therefore satisfies
    \begin{equation}
      A_n^{\max}
      =
      \frac{c_\chi}{\sqrt{\lambda_n}} .
    \end{equation}

  \subsubsection*{Exact admissibility envelope}

    We thus obtain the non-perturbative result
    \begin{equation}
      A_n^{\max}
      =
      \frac{c_\chi}{\sqrt{\lambda_n}},
    \end{equation}
    which coincides exactly with the admissibility envelope derived from the
    linear bounded-flux argument.

    The weakly non-linear correction therefore does not modify the
    location of the admissibility boundary.
    Instead, the DBI dynamics smoothly interpolates between the linear regime
    and a velocity-saturated regime as the amplitude approaches the bound.

    This shows that the spectral admissibility envelope identified above is not
    merely a linear approximation but the exact saturation boundary for the
    single-mode DBI dynamics on the $Q_8$ Cayley graph.
